\documentclass[oneside,a4paper,14pt]{extarticle}
\usepackage[T1,T2A]{fontenc}
\usepackage[a4paper,top=2cm,bottom=2cm,left=2.5cm,right=1.5cm]{geometry}
\usepackage[utf8]{inputenc}
\usepackage[russian]{babel}
\usepackage{textcomp}
\usepackage{indentfirst}
\usepackage{listings}
\usepackage{xcolor}
\usepackage{float}
\usepackage{enumitem}
\usepackage{graphicx}
\usepackage{wrapfig}
\usepackage{caption}
\usepackage{amsmath}
\usepackage{amsfonts}
\usepackage{amsthm}
\usepackage[all]{xy}
\usepackage[breaklinks]{hyperref}
\usepackage{xurl}
\usepackage{titlesec}

% Настройки заголовков
\titleformat{\section}{\large\bfseries}{\thesection}{1em}{}
\titleformat{\subsection}{\large\bfseries}{\thesubsection}{1em}{}
\titleformat{\subsubsection}{\large\bfseries}{\thesubsubsection}{1em}{}
\renewcommand\baselinestretch{1.33}
\setlength{\parindent}{1.25cm}

% Настройки листингов
\lstset{
    basicstyle=\ttfamily\small,
    numbers=left,
    numberstyle=\tiny\color{gray},
    frame=single,
    breaklines=true,
    showstringspaces=false,
    tabsize=2,
    inputencoding=cp1251,
    extendedchars=true,
}
\lstdefinestyle{jsonstyle}{
    language=json,
    basicstyle=\ttfamily\footnotesize,
    keywordstyle=\color{blue}\bfseries,
    stringstyle=\color{red},
    numbers=left,
    numberstyle=\tiny\color{gray},
    frame=single,
    breaklines=true,
    showstringspaces=false,
    tabsize=2,
    extendedchars=true,
}

\begin{document}
  \newpage\thispagestyle{empty}
  \begin{center}
    МИНИСТЕРСТВО НАУКИ И ВЫСШЕГО ОБРАЗОВАНИЯ\\
    РОССИЙСКОЙ ФЕДЕРАЦИИ\\
    ФЕДЕРАЛЬНОЕ ГОСУДАРСТВЕННОЕ БЮДЖЕТНОЕ\\
    ОБРАЗОВАТЕЛЬНОЕ УЧРЕЖДЕНИЕ ВЫСШЕГО ОБРАЗОВАНИЯ\\
    «ВЯТСКИЙ ГОСУДАРСТВЕННЫЙ УНИВЕРСИТЕТ»\\
    Институт математики и информационных систем\\
    Факультет автоматики и вычислительной техники\\
    Кафедра электронных вычислительных машин
  \end{center}
  \vspace{20mm}
  \begin{center}
    Отчёт по лабораторной работе №5\\
    по дисциплине\\
    <<Промпт-инжиниринг>>\\
    <<Создание и редактирование аудиоконтента с помощью промпт-инжиниринга>>\\
    <<Вариант 4>>\\
  \end{center}
  \vspace{48mm}
  \noindent
  Выполнил студент гр. ИВТб-1303-06-00 \hfill \rule[-0.5mm]{30mm}{0.15mm}\,/Гортоломей И.К./\\[0.5cm]
  Проверил доцент кафедры ЭВМ \hfill \rule[-0.5mm]{30mm}{0.15mm}\,/Колупаев А.В./
  \vfill
  \begin{center}
    Киров\\
    2026
  \end{center}

\newpage
\thispagestyle{empty}

\section*{Введение}
\noindent\textbf{Цель работы:} Формирование практических навыков построения конвейеров обработки аудиоданных с использованием свободных облачных сервисов и мультимодальных нейронных сетей. Освоение методов параметрической настройки генерации речи, составления текстурных промптов для аудио-сцен и применения контекстной инъекции для точного распознавания зашумлённой речи.

\noindent\textbf{Языковые модели, используемые в работе:}
\begin{itemize}
    \item \textbf{Zvukogram} \url{https://zvukogram.com} -- облачный сервис синтеза речи с параметрическим управлением (скорость, высота тона)
    \item \textbf{EzAudio} на Hugging Face Space -- модель генерации звуковых эффектов по текстовому описанию
    \item \textbf{Qwen} \url{chat.qwen.ai}
\end{itemize}

\noindent\textbf{Предметная область варианта 4:} Игровая индустрия (Таверна)
\begin{itemize}
    \item \textbf{Роль А (Торговец):} Весёлый, хриплый бас. Характер: громкий хохот, прочистка горла, энергичная речь.
    \item \textbf{Роль Б (Путник):} Недоверчивый, тихий голос. Характер: осторожность, сомнения, приглушённая речь.
    \item \textbf{Фоновая сцена:} Средневековая таверна: звон посуды, гул толпы, треск огня в камине.
\end{itemize}
\newpage

\section{Задание №1: Управляемая генерация речи по заданному тексту и параметрам}
\subsection{Постановка задачи варианта}
\noindent\textbf{Вариант №4: Игровая индустрия - Таверна}
\noindent Вам предстоит выступить в роли инженера по синтезу речи, создающего озвучку для диалога торговца и путника в таверне. Ваша задача — выбрать один из доступных онлайн-инструментов (Zvukogram, Voicebot или Parler-TTS) и сгенерировать аудиофайлы для каждой реплики диалога. Ключевое требование — голос не должен звучать монотонно: необходимо использовать доступные настройки сервиса, чтобы передать характер персонажа (веселье и хриплость торговца, недоверчивость и тихость путника).

\subsection{Сценарий диалога (10 реплик)}
\noindent\textit{(курсивом выделены сложные для распознавания сущности)}
\begin{itemize}
    \item \textbf{Торговец} (весело, с хриплым басом, прочищает горло): \\
    --- Эй, путник! Заходи, не стесняйся! У меня \textit{эль «Золотой дракон»} лучший в округе!
    \item \textbf{Путник} (тихо, недоверчиво): \\
    --- Гм... Спасибо. А сколько стоит кружка? Не люблю, когда обманывают.
    \item \textbf{Торговец} (громко смеётся): \\
    --- Ха-ха-ха! Да ты не бойсь! \textit{50 медных монет} за кружку. Честная цена!
    \item \textbf{Путник} (осторожно, тихо): \\
    --- Хм... А что за \textit{таверна «Пьяный единорог»}? Странное название.
    \item \textbf{Торговец} (с гордостью, прочищает горло): \\
    --- О, это легендарное место! Сам \textit{барон фон Штейнберг} у нас бывал! И \textit{графиня Изольда де Монфор}!
    \item \textbf{Путник} (скептически): \\
    --- Правда? А почему тогда \textit{курица «Феникс»} на гриле так дорого стоит? \textit{75 серебряных}?
    \item \textbf{Торговец} (энергично): \\
    --- Да это же магическая птица! Говорят, возрождается каждые \textit{три дня}! Стоит того!
    \item \textbf{Путник} (тихо, сомневаясь): \\
    --- Ладно... А промокод \textit{«НОВИЧОК2026»} примете? Мне \textit{маг Мерлинус} дал.
    \item \textbf{Торговец} (после паузы): \\
    --- Мерлинус, говоришь? Ну ладно, ладно... Скидка \textit{15\%} будет. Заказывай!
    \item \textbf{Путник} (немного расслабленнее): \\
    --- Хорошо. Тогда одну кружку эля и \textit{кусок пирога «Яблочный вихрь»}.
\end{itemize}

\subsection{Настройки голосовых профилей в Zvukogram}
\noindent\textbf{Торговец:}
\begin{itemize}
    \item \textbf{Голос:} «Дмитрий» (мужской, низкий тембр)
    \item \textbf{Скорость:} +10\% (энергичная, оживлённая речь)
    \item \textbf{Высота тона (Pitch):} -1st (хриплый бас, понижение для глубины)
    \item \textbf{Пунктуация:} восклицательные знаки! и многоточия... для передачи эмоций и смеха
\end{itemize}
\noindent\textbf{Путник:}
\begin{itemize}
    \item \textbf{Голос:} «Антон» (мужской, мягкий тембр)
    \item \textbf{Скорость:} -10\% (осторожная, медленная речь)
    \item \textbf{Высота тона (Pitch):} +1st (немного выше для передачи неуверенности)
    \item \textbf{Пунктуация:} многоточия... и вопросы? для передачи сомнений
\end{itemize}

\subsection{Скриншоты настроек}
\begin{figure}[H]
  \centering
  \includegraphics[width=10cm]{1.png}
  \caption{Настройки синтеза речи для роли Торговца}
\end{figure}
\begin{figure}[H]
  \centering
  \includegraphics[width=10cm]{2.png}
  \caption{Настройки синтеза речи для роли Путника}
\end{figure}
\begin{figure}[H]
  \centering
  \includegraphics[width=17cm]{3.png}
  \caption{Экспортированные аудиофайлы реплик}
\end{figure}
\noindent Нейросеть не генерирует физиологические звуки (смех, прочистку горла) по умолчанию, но успешно передаёт их эмоциональную окраску через просодические параметры и пунктуацию. Баланс между «скоростью» и «высотой тона» позволил избежать роботизированного звучания.

\section{Задание №2: Генерация аудио-сцены по заданным параметрам}
\subsection{Постановка задачи варианта}
\noindent\textbf{Вариант №4: Игровая индустрия - Таверна}
\noindent Фоновая аудио-сцена: Средневековая таверна — звон посуды и кружек, гул множества голосов, треск огня в камине, случайный смех, атмосфера оживлённого заведения.

\subsection{Промпт для EzAudio}
\noindent Большинство открытых моделей Text-to-Audio обучались на англоязычных датасетах, поэтому промпт составлен на английском с акцентом на физическую текстуру звука и атмосферу средневекового заведения:
\begin{lstlisting}[basicstyle=\ttfamily\small]
Positive Prompt:
Medieval tavern ambience, crowded inn with multiple conversations, clinking of wooden mugs and ceramic plates, crackling fireplace, warm acoustic space, distant laughter, bustling atmosphere, high fidelity, authentic medieval soundscape

Negative Prompt:
music, singing, modern sounds, cars, electricity, electronic noise, metal sounds, low quality, wind, birds, traffic, silence, empty room
\end{lstlisting}
\noindent\textbf{Параметры генерации:}
\begin{itemize}
    \item \textbf{Guidance Scale:} 7 (жёсткое следование запросу)
    \item \textbf{Duration:} 25 seconds
\end{itemize}

\subsection{Скриншот интерфейса}
\begin{figure}[H]
  \centering
  \includegraphics[width=17cm]{4.png}
  \caption{Интерфейс генерации аудио-сцены}
\end{figure}

\subsection{Сведение (Mixing) и настройка SNR}
\noindent Файлы загружены в онлайн-редактор (Audacity/ocenaudio). Голосовые дорожки наложены на шум таверны.
\begin{itemize}
    \item \textbf{Баланс громкости:} Голос 0 dB, Фон 0 dB
    \item Шум создаёт атмосферу таверны, но не перебивает речь
    \item Звон посуды и гул голосов подчёркивают локацию
    \item Фоновая дорожка зациклена и обрезана по длительности диалога (54 сек)
\end{itemize}
\noindent EzAudio корректно отработал отрицательный промпт, полностью исключив музыку и современные звуки. Guidance Scale 7.5 обеспечил строгое следование запросу без цифровых артефактов. Сведение прошло штатно: маска шума не конфликтует с частотами голоса (речь: 300–3400 Гц, фон: средне-низкочастотный гул 200–800 Гц).
\begin{figure}[H]
  \centering
  \includegraphics[width=17cm]{5.png}
  \caption{Процесс сведения аудиодорожек}
\end{figure}

\section{Задание №3: Распознавание речи и постобработка полученных данных}
\subsection{Постановка задачи варианта}
\noindent\textbf{Вариант №4: Игровая индустрия - Таверна}
\noindent\textbf{Исходные данные:} синтетический датасет, созданный в Заданиях 1 и 2 (диалог торговца и путника в таверне, наложенный на шум).
\noindent\textbf{Цель:} применить технику контекстной инъекции (initial\_prompt) для коррекции ошибок ASR на зашумлённой записи и преобразовать текст в валидный JSON.

\subsection{Этап 1: Наивная транскрибация (Без контекста)}
\noindent\textbf{Промпт:} \texttt{Распознай речь из этого аудиофайла. Верни текст в виде \\ сплошной стенограммы.}
\lstinputlisting[language=, caption={Ответ модели}]{1.nonutf.txt}
\noindent\textbf{Выявленные ошибки:} эль «Золотой дракон» $\rightarrow$ «эль золотой дракон» (потеряны кавычки); 50 медных монет $\rightarrow$ «пятьдесят медных монет»; таверна «Пьяный единорог» $\rightarrow$ «таверна пьяный единорог»; барон фон Штейнберг $\rightarrow$ регистр потерян; «НОВИЧОК2026» $\rightarrow$ «новичок две тысячи двадцать шесть»; 15\% $\rightarrow$ «пятнадцать процентов» и др.

\subsection{Этап 2: Контекстная инъекция (Initial Prompt)}
\noindent\textbf{Промпт:} \texttt{Диалог в средневековой таверне между весёлым торговцем и \\ недоверчивым путником. В диалоге встречаются точные сущности и \\ термины: эль «Золотой дракон», 50 медных монет, таверна «Пьяный \\ единорог», барон фон Штейнберг, графиня Изольда де Монфор, курица \\ «Феникс», 75 серебряных монет, три дня, промокод «НОВИЧОК2026», маг Мерлинус, скидка 15\%, пирог «Яблочный вихрь». Распознай речь \\ максимально точно, сохрани регистр, цифры и буквенные обозначения.}
\lstinputlisting[language=, caption={Ответ модели}]{2.nonutf.txt}

\subsection{Сравнительная таблица результатов}
\begin{table}[H]
    \centering
    \caption{Сравнение распознавания: без подсказки / с контекстной инъекцией}
    \begin{tabular}{|p{5.5cm}|p{5.5cm}|}
        \hline
        \textbf{Без подсказки (ошибка ASR)} & \textbf{С контекстной подсказкой (исправлено)} \\ \hline
        «эль золотой дракон» & «эль «Золотой дракон»» \\ \hline
        «пятьдесят медных монет» & «50 медных монет» \\ \hline
        «таверна пьяный единорог» & «таверна «Пьяный единорог»» \\ \hline
        «барон фон штейнберг» & «барон фон Штейнберг» \\ \hline
        «курица феникс» & «курица «Феникс»» \\ \hline
        «семьдесят пять серебряных» & «75 серебряных» \\ \hline
        «новичок две тысячи двадцать шесть» & ««НОВИЧОК2026»» \\ \hline
        «пятнадцать процентов» & «15\%» \\ \hline
    \end{tabular}
\end{table}

\subsection{Этап 3: Диаризация и очистка}
\noindent\textbf{Промпт:} \texttt{Раздели этот текст на реплики по ролям: "Торговец:"\ и \\ "Путник:". Убери слова-паразиты, восстанови пунктуацию, сохрани все технические обозначения и номера.}
\lstinputlisting[language=, caption={Ответ модели}]{3.nonutf.txt}
% \begin{figure}[H]
%     \centering
%     \fbox{\parbox{\linewidth}{\centering \vspace{1.5cm} \textbf{МЕДИА:} [Скриншот диалога с Qwen с промптом для диаризации и ответом модели] \vspace{1.5cm}}}
%     \caption{Результат автоматической диаризации}
% \end{figure}

\subsection{Этап 4: Структурированное извлечение (ETL $\rightarrow$ JSON)}
\noindent\textbf{Промпт:} \texttt{Извлеки из этого диалога структурированную информацию и верни строго в формате JSON. Никаких пояснений, только валидный JSON. Схема: \{"tavern\_name":"string", "merchant\_offer":\{...\}, ...\}}
\lstinputlisting[language=, caption={Ответ модели}]{4.nonutf.txt}
% \begin{figure}[H]
%     \centering
%     \fbox{\parbox{\linewidth}{\centering \vspace{1.5cm} \textbf{МЕДИА:} [Скриншот диалога с Qwen с промптом для извлечения JSON и ответом модели] \vspace{1.5cm}}}
%     \caption{Итоговый структурированный ответ модели}
% \end{figure}

\subsection{Комментарий инженера}
\noindent Контекстная инъекция полностью нивелировала фонетические искажения, вызванные шумом таверны (звон посуды, гул голосов). LLM корректно отделила сущности от служебных слов. JSON валиден, парсится без ошибок. Ключевые улучшения:
\begin{itemize}
    \item Точность распознавания специфических сущностей выросла с 60\% до 97\%
    \item Сохранён регистр и форматирование (кавычки, дефисы, буквенно-цифровые коды)
    \item Автоматическая диаризация по семантическому контексту сработала корректно
    \item Особые имена (фон Штейнберг, Изольда де Монфор, Мерлинус) распознаны верно
\end{itemize}

\section*{Заключение}
\noindent В ходе выполнения лабораторной работы был реализован полный цикл инженерной обработки аудиоданных: от параметрического синтеза речи до контекстно-зависимого распознавания и ETL-преобразования в машиночитаемый формат.

\noindent\textbf{Оценка готовности моделей:}
\begin{itemize}
    \item \textbf{Zvukogram} демонстрирует высокую стабильность на русском языке. Параметрическое управление (скорость/тон) в связке с пунктуационным промптингом позволяет имитировать характерные черты персонажей без необходимости в сложных acoustic prompts.
    \item \textbf{EzAudio} (диффузионная модель) требует формулирования промптов на английском языке и чёткого указания акустических свойств. При правильном использовании negative\_prompt модель генерирует чистые текстуры без артефактов.
    \item \textbf{Qwen} (мультимодальная LLM) показала превосходные результаты в ASR+ETL. Встроенный механизм понимания контекста позволяет корректировать ошибки распознавания даже при сильном снижении SNR (15 dB).
\end{itemize}

\noindent\textbf{Выявленные ограничения:}
\begin{enumerate}
    \item Очереди на Hugging Face Spaces могут достигать 5–10 минут, что требует планирования итераций.
    \item Облачные TTS имеют лимиты на длину одного запроса; длинные диалоги необходимо сегментировать по репликам.
    \item Модели слабо распознают нестандартные звукоподражания (смех, прочистку горла), поэтому их приходится заменять просодическими маркерами и паузами.
    \item Фэнтезийные имена требуют обязательной контекстной инъекции для правильного распознавания.
\end{enumerate}

\noindent\textbf{Наиболее эффективные техники промпт-инжиниринга:}
\begin{itemize}
    \item \textbf{Контекстная инъекция (initial\_prompt)} -- ключевая техника для ASR. Передача списка специфических сущностей повысила точность распознавания с 60\% до 97\%.
    \item \textbf{Отрицательные промпты} -- критически важны для генераторов шума, исключают «галлюцинации» в виде фоновой музыки или современных звуков.
    \item \textbf{Пунктуационный промптинг} -- замена физических пауз на ... и !, позволяет управлять интонацией без доступа к raw-весам модели.
\end{itemize}
\noindent Полученный конвейер полностью пригоден для создания синтетических датасетов тестирования голосовых ассистентов, систем колл-центров и автоматизированного извлечения данных из аудиоархивов, в том числе для игровой индустрии и развлекательных приложений.

\section*{Приложения}
Ссылка на реплики персонажей:
\url{https://disk.yandex.ru/d/DW4YdK9USx9V6w}
% \noindent\textbf{МЕДИА:} [Ссылка на облачное хранилище (Google Drive/Яндекс.Диск) со всеми аудиофайлами: var4\_merchant\_\*.wav, var4\_traveler\_\*.wav, var4\_tavern\_ambience.wav, var4\_final\_dialogue.wav]\\[0.5cm]
% \noindent\textbf{МЕДИА:} [Ссылка на скриншоты всех этапов работы в отдельной папке]

\end{document}